\section{ĐỀ SỐ \thesection } %đề 18
\setcounter{bt}{0}

%%Câu 1.
\begin{bt} Cho tập hợp $D=\{x\in Z|-3<x\leq 1\}$. \\
	Em hãy viết lại tập hợp $D$ bằng cách liệt kê các phần tử.
	\loigiai{$D=\{-2;-1;0;1\}$}
\end{bt}


%%Câu 2.
\begin{bt} Tính nhanh (nếu có thể):
	\begin{listEX}[2]
		\item $66\cdot(-45)+(-55)\cdot66$;
		\item $257-7[(-3)^2-(17-16)^{2023}]$;
		\item $[(3^2+2^3)\cdot2-52:13]:(5\cdot3)$;
		\item $2[(25-3^{2024}:3^{2022}):2^2+149]-400$.
	\end{listEX}
	\loigiai{
		\begin{listEX}[2]
				\item $66\cdot(-45)+(-55)\cdot66\\
				=66(-45-55)\\
				=66\cdot(-100)
				=-6600$
				\item $257-7[(-3)^2-(17-16)^{2023}]\\=257-7[9-1]\\=257-7\cdot8=257-56=201$
				\item $[(3^2+2^3)\cdot2-52:13]:(5\cdot3)\\=[(9+8)\cdot2-4]:15\\=(17\cdot2-4)\cdot15\\=30:15\\=2$
				\item $2[(25-3^{2024}:3^{2022}):2^2+149]-400\\=2[(25-3^2):4+149]-400\\=2\cdot (16:4+149)-400\\=2\cdot153-400=-94$
			\end{listEX}
	}
\end{bt}



%%Câu 3.
\begin{bt} Tìm số nguyên $x$, biết:
	\begin{listEX}[2]
		\item $x-42=-8$;
		\item $(x+12):15=60$;
		\item $36+(100-5x):4=51$;
		\item $5(x+1)^2-21=-1$.
	\end{listEX}
	\loigiai{\begin{multicols}{2}
			\begin{enumerate}
				\item 
				$\begin{aligned}
					x - 42 &= -8\\
					x &= -8 + 42 = 34
				\end{aligned}$
				
				\item 
				$\begin{aligned}
					36 + (100 - 5x):4 &= 51\\
					(100 - 5x):4 &= 51 - 36 = 15\\
					100 - 5x &= 15\cdot4 = 60\\
					5x &= 100 - 60 = 40\\
					x &= 40:5 = 8
				\end{aligned}$
				
				\item 
				$\begin{aligned}
					(x + 12):15 &= 60\\
					x + 12 &= 60\cdot15 = 900\\
					x &= 900 - 12 = 888
				\end{aligned}$
				
				\item 
				$\begin{aligned}
					5(x + 1)^2 - 21 &= -1\\
					5(x + 1)^2 &= -1 + 21 = 20\\
					(x + 1)^2 &= 20:5 = 4
				\end{aligned}$\\
				$\Rightarrow x + 1 = 2$ hoặc $x + 1 = -2$\\
				$\Rightarrow x = 1$ hoặc $x = -3$
			\end{enumerate}
		\end{multicols}
		
	}
\end{bt}

%%Câu 4.
\begin{bt} An, Bình, Chi cùng học một đường. An cứ 5 ngày trực nhật một lần, Bình 10 ngày và Chi 8 ngày một lần. Lần đầu cả ba bạn cùng trực nhật vào một hôm. Hỏi sau ít nhất bao nhiêu ngày thì ba bạn lại cùng trực một hôm?
	\loigiai{Gọi $x$ $(x\in  \mathbb{N})$ là số ngày ít nhất để ba bạn lại cùng trực một hôm.Ta có $x=BCNN(5,10,8)$\\
		 Phân tích mõi số ra thừa số nguyên tố 
			$5=5$ \hspace{1cm} $10=2\cdot5$\hspace{1cm} $8=2^3$
		 
		$BCNN(5,10,8)=2^3\cdot5=40$.\\
		Vậy sau 40 ngày ba bạn lại cùng trực nhật một hôm.
	}
\end{bt}


%%Câu 5.

\begin{bt} \immini[thm]{
		Bác Bảo có một khu đất hình vuông có kích thước 35 m.
		\begin{enumerate}
			\item Tính chu vi và diện tích khu đất của bác Bảo.
			\item Bác Bảo chia mảnh đất làm hai như hình bên. Bác cần rào phần đất nuôi bò lại để làm chuồng. Hỏi bác Bảo cần bao nhiêu mét kẽm gai để vừa đủ rào mảnh đất làm chuồng bò? Tính diện tích chuồng bò.
		\end{enumerate}
	}{
		\begin{tikzpicture}[scale=1,line width=0.8pt,line join=round,>=stealth,font=\small]
			\def\s{9}
			;
			\draw[fill=white] (0,0) rectangle (35/\s,35/\s);
			\draw[fill=lightgray] (0,0) --(0,10/\s)--(35/\s,18/\s)node[midway,sloped,above]{36 m}--(35/\s,0)--cycle;
			\draw (0.5,30/\s) node[right]{Đất trồng cỏ} (35/\s-0.5,0.5)node[left]{Chuồng bò};
			\draw[<->] (-0.3,10/\s)--++(90:25/\s) node[midway,above, sloped]{25 m};
			\draw[<->] (35/\s+0.3,18/\s)--++(90:17/\s) node[midway,below,sloped]{17 m};
		\end{tikzpicture}
	}
	\loigiai{\begin{enumerate}
			\item Chu vi khu đất: $35\cdot4=140$ (m)\\Diện tích khu đất: $35\cdot35=1225$ (m$^2$)
			\item Số mét kẽm cần: $10+18+35+36=99$ (m) \\ Diện tích chuồng bò: $(35-25+35-17)\cdot35:2=490$ (m$^2$)
	\end{enumerate}}
\end{bt}



%%Câu 6.
\begin{bt} Vẽ Biểu đồ cột biểu diễn môn thể thao yêu thích của học sinh lớp 6A được cho trong bảng thống kê sau:
	\begin{center}
		\begin{tabular}{|c|c|c|c|c|}\hline
			\textbf{Môn thể thao}&\textbf{Bóng đá}&\textbf{Cầu lông}&\textbf{Bóng bàn}&\textbf{Bóng rổ} \\ \hline
			\textbf{Số học sinh yêu thích}&10&7&3&5 \\ \hline
		\end{tabular}\\
	\end{center}
	\loigiai{
		\begin{center}
			\begin{tikzpicture}
				[scale=0.8,font=\small,line join=round,>=stealth]
				\def\hoanh{10.2};
				\def\tung{4.2};
				\node[above] at (\hoanh/2,\tung+1){\normalsize \bfseries Môn thể thao yêu thích của học sinh lớp 6A};
				\def\mau{mauChuDao};
				\draw (0,0)node[left]{0}--(\hoanh,0);
				\draw[->] (0,0)--(0,\tung+1);
				\foreach \y in {2,4,...,12}{\draw (-0.1,\y/3)node[left]{$\y$}--++(0:0.2);
					\draw[dashed,gray] (0,\y/3)--(\hoanh-0.2,\y/3);}
				\foreach \x/\n in{2/10,4/7,6/3,8/5}{\draw[line width=7mm,\mau] (\x,0)--++(0,\n/3)node[black,shift={(90:0.3)}]{\n};}
				\foreach \x/\p in {2/Bóng đá,4/Cầu lông,6/Bóng bàn,8/Bóng rổ}{\node at (\x,-0.4){\footnotesize \p};} 
			\end{tikzpicture}
		\end{center}
	}
\end{bt}


%%Câu 7.
\begin{bt} Anh Long có tiền lương cố định hàng tháng là $8\,000\,000$ đồng. Nếu anh Long làm tăng ca thì được tính thêm $200\,000$ đồng một giờ và nếu anh làm sai kỹ thuật thì 1 lỗi sẽ bị trừ $150\,000$ đồng. Tháng 12, anh Long làm tăng ca 28 giờ và bị sai 4 lỗi kỹ thuật. Hỏi tiền lương thực lãnh của anh Long trong tháng 12 là bao nhiêu?
	
	\loigiai{Tiền lương của anh Long tháng 12 là:\\ $8\,000\,000+28\cdot200\,000-4\cdot150\,000=13\,000\,000$ (đồng)}
\end{bt}



